\section{Ответы на контрольные вопросы}
\label{sec:questions}

\subsection{Вопрос 1. Зависимость числа итераций от коэффициента обучения}

В варианте 2 коэффициент обучения не является постоянным: для каждого образа он вычисляется по формуле~\eqref{eq:alpha}. Свободным параметром, который задаёт пользователь, является масштаб $\beta$. Поэтому исследуется зависимость $T(\beta)$.

В окрестности минимума функцию потерь можно аппроксимировать квадратичной формой
\begin{equation}
	E(w) \approx \frac{1}{2}(w-w^*)^{\mathsf{T}}H(w-w^*),
\end{equation}
где $H=H^{\mathsf{T}}$~--- матрица Гессе. Обозначим $u(t)=w(t)-w^*$. Градиентный шаг даёт
\begin{equation}
	u(t+1)=\bigl(I-\alpha H\bigr)u(t).
\end{equation}
По спектральной теореме $H=V\Lambda V^{\mathsf{T}}$, где $V$ ортогональна, $\Lambda=\operatorname{diag}(\lambda_1,\ldots,\lambda_n)$. В базисе $z=V^{\mathsf{T}}u$ система распадается на независимые компоненты:
\begin{equation}
	z_i(t+1)=(1-\alpha\lambda_i)z_i(t),\qquad i=1,\ldots,n.
	\label{eq:q1-eigen}
\end{equation}
Каждая компонента затухает как геометрическая прогрессия со знаменателем $q_i=1-\alpha\lambda_i$. Условие сходимости $|q_i|<1$ даёт
\begin{equation}
	0<\alpha<\frac{2}{\lambda_{\max}}.
	\label{eq:q1-stab}
\end{equation}
Число итераций определяется наиболее медленно затухающей компонентой. При $0<\alpha<1/\lambda_{\max}$ увеличение шага ускоряет затухание; при $1/\lambda_{\max}<\alpha<2/\lambda_{\max}$ компонента с $\lambda_{\max}$ начинает колебаться, и число итераций растёт. Поэтому зависимость $T(\beta)$ имеет минимум вблизи $\beta\approx1/\lambda_{\max}$, а за границей~\eqref{eq:q1-stab} метод расходится.

Экспериментальные данные (рисунок~\ref{fig:family-a}, таблица~\ref{tab:family-a}) качественно подтверждают этот вывод. При малых $\beta$ число эпох велико: для контуров при $\beta=0{,}1$ требуется 39 эпох, а при $\beta=0{,}5$~--- уже 9. Минимум $T$ достигается при $\beta=0{,}5$ для плавных переходов и текстуры и при $\beta=0{,}75$ для контуров. При $\beta\ge1{,}25$ контуры и текстура перестают достигать порога, а плавные переходы расходятся уже при $\beta=2$. Для разных изображений критическое $\beta$ различается из-за разного спектра $X^{\mathsf{T}}X$.

\subsection{Вопрос 2. Зависимость числа итераций от коэффициента сжатия}

Рассмотрим матрицу вторых моментов выборки
\begin{equation}
	R=\frac{1}{Q}X^{\mathsf{T}}X,
	\label{eq:q2-r}
\end{equation}
где $X$ имеет размер $Q \times n$. Матрица $R$ симметрична и положительно полуопределена:
\begin{equation}
	v^{\mathsf{T}}Rv=\frac{1}{Q}\|Xv\|^2\ge0.
\end{equation}
Её собственные числа $\mu_1\ge\mu_2\ge\ldots\ge\mu_n\ge0$ описывают энергию независимых направлений данных.

Для линейного сжатия размерности $n\to p$ наилучшее приближение по среднеквадратичной ошибке достигается проектированием на первые $p$ собственных векторов $R$~\cite{golovko}. Минимальная достижимая ошибка равна сумме отброшенных собственных чисел:
\begin{equation}
	E_{\min}(p)=\sum_{i=p+1}^{n}\mu_i.
	\label{eq:q2-emin}
\end{equation}
При фиксированных $Q$ и $n$ коэффициент $K$ убывает с ростом $p$ (приближённо $K\sim1/p$). Уменьшение $p$ увеличивает $K$, но одновременно увеличивает $E_{\min}$: отбрасывается больше направлений.

Если $E_{\min}(p)\ge\varepsilon$, то порог $\varepsilon$ недостижим для любого линейного преобразования ранга не выше $p$, независимо от числа эпох. Если $E_{\min}(p)<\varepsilon$, то зазор $\varepsilon-E_{\min}$ определяет, насколько быстро ошибка может пересечь порог: чем меньше зазор, тем больше эпох требуется. Этим объясняется появление прочерков в таблице~\ref{tab:family-b}: для контуров при $p=12$ ($K=1{,}684$) порог не достигается, а при $p=16$ ($K=1{,}263$)~--- достигается за 12 эпох; для текстуры аналогично $p=12$ недостижим, а $p=16$ даёт 6 эпох. Максимальное достижимое сжатие ограничено спектром конкретной выборки.

\subsection{Вопрос 3. Зависимость числа итераций от вида изображения}

Вид изображения определяет матрицу $R$ и её спектр. Для изображений с высокой пространственной корреляцией (плавные градиенты, крупные однородные области) соседние пиксели изменяются согласованно, и большая часть энергии сосредоточена в нескольких первых главных компонентах: $\mu_1\gg\mu_2\gg\ldots$ При фиксированном $p$ сумма отброшенных чисел~\eqref{eq:q2-emin} мала, зазор $\varepsilon-E_{\min}$ велик, и число эпох невелико.

Для изображений с резкими границами и мелкой текстурой корреляция между соседними пикселями слабее, энергия распределена по большему числу направлений, и хвост спектра убывает медленно. При том же $p$ величина $E_{\min}$ больше, зазор меньше, поэтому требуется больше эпох или порог вообще недостижим при данном $p$.

Экспериментальные данные (рисунок~\ref{fig:family-c}, таблица~\ref{tab:family-c}) согласуются с этим при одинаковом требовании к ошибке: при общей $\mathrm{MSE}_n=\CommonMSE$ плавные переходы достигают порога за 1 эпоху, текстура~--- за 2, контуры~--- за 4. При сравнении изображений важно фиксировать одинаковое требование к ошибке: при разных $\varepsilon$ число эпох отражает не только сложность изображения, но и выбранный порог. Поэтому в таблице~\ref{tab:family-c} приведены оба варианта~--- индивидуальные пороги и общая допустимая MSE.

Индивидуальные пороги выбраны по ошибке 10-й эпохи, и для текстуры траектория ошибки немонотонна: после 6-й эпохи ошибка снова возрастает, поэтому первое пересечение индивидуального порога происходит раньше ($T=6$), чем у плавных переходов ($T=9$). Это не противоречит выводу: при общем пороге, который не зависит от случайных колебаний отдельной траектории, упорядочение изображений по сложности сохраняется.

\subsection{Вопрос 4. Зависимость числа итераций от порога ошибки}

Из соотношения~\eqref{eq:q1-eigen} для наиболее медленно затухающей компоненты при малых $\alpha$ имеем
\begin{equation}
	z_{\min}(t)\approx(1-\alpha\lambda_{\min})^t z_{\min}(0).
\end{equation}
Ошибка относительно минимума пропорциональна квадрату этой компоненты:
\begin{equation}
	E(t)-E^*\approx C\,(1-\alpha\lambda_{\min})^{2t},
	\label{eq:q4-decay}
\end{equation}
где $C$ зависит от начальных условий, а $E^*$ совпадает с минимально достижимой ошибкой $E_{\min}$.

Условие достижения порога $E(T)-E^*=\varepsilon-E^*$ после логарифмирования даёт
\begin{equation}
	T(\varepsilon)=
	\left\lceil
	\frac{\ln\!\bigl(C/(\varepsilon-E^*)\bigr)}
	     {-2\ln(1-\alpha\lambda_{\min})}
	\right\rceil .
	\label{eq:q4-T}
\end{equation}
Отсюда следуют два вывода. Во-первых, $T$ зависит от $\varepsilon$ логарифмически: уменьшение порога в $k$ раз увеличивает число итераций на аддитивную величину порядка $\ln k/(2\alpha\lambda_{\min})$. Во-вторых, при $\varepsilon\to E^*$ знаменатель стремится к нулю, а $T\to\infty$: порог ниже минимально достижимой ошибки не достигается ни за какое конечное число эпох.

Данные семейства D (рисунок~\ref{fig:family-d}, таблица~\ref{tab:family-d}) показывают качественно логарифмический рост $T$ при ужесточении $\varepsilon$ и резкое увеличение числа эпох вблизи предельной ошибки. Для разных изображений предельная ошибка $E^*$ различна, поэтому одна и та же величина $\varepsilon$ может быть достижима для одного изображения и недостижима для другого.
